\section{Full optimizer output for residual-convex cubic recourse}
\label{sec:recourse}

Consider an objective with a small group of core variables $v$ and many
residual variables $z$. Convex residual problems give efficient value
queries at rational cores. They do not by themselves give a point near
the residual optimizer at an unknown exact optimal core. This section
closes that gap for cubics on product boxes, using one finite perturbation
of the core coefficients. The result strengthens the earlier cubic
completion theorem by removing its supplied joint quadratic core
convexifier. The
\href{../cubic-recourse/theorem.md}{complete companion proof} and its
\href{../reviews/cubic-extension/review.md}{independent review}
record the inherited interfaces and their verification scope.

\subsection{The output contract}

Let $F\in\Q[v,z]$ have total degree at most three on
$[0,1]^k\times[0,1]^n$. Assume that $F(v,\cdot)$ is convex for every
$v$, and supply rational $L\ge0$ with $F_{v_iv_i}\le L$ throughout
the box. Let $\sigma>0$ be rational. The input length $I$ includes
these data and any structural certificates; polynomial certificate
verification is assumed, or its actual cost must be added. Fixed
coordinates may first be substituted and positive-width boxes
normalized. Perturb only the core:
\[
 F_\gamma(v,z)=F(v,z)+\gamma^Tv,
 \qquad \gamma\in[-\sigma,\sigma]^k.
\]
For each draw, let $a$ be the lexicographically first globally optimal
core. Define
\[
 S_a=\argmin_{z\in[0,1]^n}F(a,z),
 \qquad p=\argmin_{z\in S_a}\norm{z}^2.
\]
Compactness and residual convexity make $p$ unique.

\begin{theorem}[Full cubic recourse point oracle]
\label{thm:cubic-recourse}
There is a base-computable endpoint-inclusive uniform product grid with
$M$ labels per coordinate and $\log M=\poly(I)$ such that, for every
draw and every integer $q\ge0$, an evaluator returns a rational feasible
$(v_q,z_q)$ satisfying
\[
 \norm{(v_q,z_q)-(a,p)}\le2^{-q}.
\]
Expected bit work and proof/output length are
\begin{equation}
 f(k)(1+L/\sigma)^k\poly(I+q),
 \label{eq:recourse-work}
\end{equation}
where $f$ is computable and the polynomial exponent is absolute.
One random work factor controls all precision queries. The noise law
and selector are fixed before those queries. A certified gap for the
sampled objective can be returned at the same order of cost.
No joint convexifier, residual strong-convexity modulus, or unique
residual optimizer is required.
\end{theorem}

For $n=0$, this is the existing selected-core theorem. For $k=0$,
Theorem~\ref{thm:cubic-domain} supplies minimum-norm point output.
Assume $k,n\ge1$ below. We use three established interfaces: the
\href{../../research-20261002/new-direction/core-only-noise-core-oracle.md}{selected-core oracle},
polynomial-time certified convex value optimization at rational cores,
and the
\href{../../research-20261002/new-direction/polynomial-exact-fallback.md}{exact polynomial fallback}.
The first returns short rational approximations to $a$ with an
all-precision work factor of the form \eqref{eq:recourse-work}. Its mesh
requirements are lower bounds fixed from base data. The third has
cost $B_0\poly(I+b+q)$, where $B_0=2^{\poly(I)}$ is independent of
sampled coefficient length $b$ and requested precision $q$.

\subsection{A sound certificate for exact core endpoints}

Put $\bar L=\max\{1,L\}$ and fix a base rational $0<t\le1/2$.
For each $i$, query selected cores of the two auxiliary objectives
$F_{\gamma-te_i}$ and $F_{\gamma+te_i}$ to coordinate error
$r\le t/(8\bar L)$. Let $[\ell_i^-,u_i^-]$ and
$[\ell_i^+,u_i^+]$ be the resulting certified intervals. Use
\begin{equation}
 u_i^-<t/\bar L\ \Longrightarrow\ v_i=0,
 \qquad
 \ell_i^+>1-t/\bar L\ \Longrightarrow\ v_i=1.
 \label{eq:recourse-face-test}
\end{equation}
Each conclusion holds for every original global optimizer, including
on tied draws.

To prove this, hold $\gamma_{-i}$ fixed and project out all other
variables:
\[
 W(s)=\min_{v_{-i},z}
       \{F(v,z)+\gamma_{-i}^Tv_{-i}:v_i=s\}.
\]
The function $W(s)-\bar Ls^2/2$ is concave, since each function in
the minimum has this property. If $a_i$ minimizes
$W(s)+\gamma_i s$ and $b_i$ minimizes
$W(s)+(\gamma_i-t)s$, their optimality inequalities imply
$a_i\le b_i$. At an interior minimizing contact, semiconcavity gives
differentiability, with derivative equal to minus the tilt. Thus, if
both coordinates are interior,
\[
 t=W'(b_i)-W'(a_i)\le\bar L(b_i-a_i).
\]
The first strict interval test contradicts this inequality whenever
$a_i>0$. Reflection proves the other test.

Failure to recognize an actual endpoint has a short conditional
noise interval. Set
\[
 \theta=\sup_{s>0}\frac{W(0)-W(s)}s.
\]
Coefficient bounds make $W$ Lipschitz, so $\theta$ is finite. An
original zero coordinate requires $\gamma_i\ge\theta$. If
$\gamma_i>\theta+t$, every negatively shifted optimizer has coordinate
zero, and the first test passes. A missed lower endpoint therefore
requires $\gamma_i\in[\theta,\theta+t]$; the upper endpoint has the
same bound. On an endpoint-inclusive $M$-label grid, an interval of
length $t$ has mass at most $t/(2\sigma)+2/M$. Consequently
\begin{equation}
 \Pr\{\text{some endpoint of the selected core is missed}\}
 \le kt/\sigma+4k/M.
 \label{eq:recourse-face-tail}
\end{equation}
Unforced coordinates are provisionally free. A later certificate rejects
any that are actually endpoints; probability is not used to justify a
face declaration.

\subsection{Cubic fibers and a lattice certificate for their minors}

Let $A$ be the certified face, with $m$ free coordinates and relative
center $c_A$, and let $c_z=(1/2,\ldots,1/2)$. Define
\[
 H_A=\nabla^2_{zz}F(c_A,c_z),\qquad
 g_A(v)=\nabla_zF(v,c_z),\qquad
 B_A(v)=\begin{bmatrix}H_A\\g_A(v)^T\\I_n\end{bmatrix}.
\]
Write $H(v,z)=\nabla^2_{zz}F(v,z)$. This matrix is affine and positive
semidefinite. Reflection about $(c_A,c_z)$ gives
$\ker H_A\subseteq\ker H(v,z)$ on the whole face product, and
$\ker H(v,c_z)=\ker H_A$ for $v\in\operatorname{relint}A$.
Individual residual boundary points can have larger kernels.
Provided $a\in\operatorname{relint}A$, as acceptance below certifies,
the cubic argument behind Theorem~\ref{thm:cubic-domain} gives, for
any $y\in S_a$,
\begin{equation}
 S_a=\{z\in[0,1]^n:H_A(z-y)=0,\ g_A(a)^T(z-y)=0\}.
 \label{eq:recourse-fiber}
\end{equation}
Only the row $g_A(v)^T$ varies with $v$, and it has degree at most two.
Every square minor of $B_A(v)$ therefore has degree at most two.

Compute integers $D,B\ge1$, with polynomial bit length, such that
every minor on every original face, multiplied by $D$, is an integer
polynomial of coefficient height at most $B$. This requires only
denominator and coefficient bounds. For example, let $d$ clear all
possible entry coefficients, and let $C$ bound their cleared heights.
With $s_*=(k+1)(k+2)/2$, the conservative choices
\[
 D=d^n,\qquad B=D\,n!(s_*C)^n
\]
suffice. Neither minors nor faces are enumerated.

For the free coordinates of $a$, let $\phi(a)$ list all monomials of
degree at most two, including one; its length is
$s=(m+1)(m+2)/2$. Choose an integer scale $T$ and obtain a short dyadic
$b$ with $\norm{b-a}_\infty\le1/(4T)$. Put
$w_j=\operatorname{round}(T\phi_j(b))$. Since these monomials are
$2$-Lipschitz on the cube,
$|w_j/T-\phi_j(a)|\le1/T$.
Apply LLL reduction with parameter $3/4$
\citep{LenstraLenstraLovasz1982} to the rank-$s$
integer lattice generated by $[I_s\mid w]$. If $\ell_1$ is its first
reduced vector, accept only when
\begin{equation}
 \norm{\ell_1}^2>2^{s-1}(4sB)^2.
 \label{eq:recourse-lattice-test}
\end{equation}
The usual inequality
$\norm{\ell_1}\le2^{(s-1)/2}\lambda_1$ then proves
\begin{equation}
 |h^T\phi(a)|>sB/T
 \quad(0<\norm{h}_\infty\le B, h\in\Z^s).
 \label{eq:recourse-all-margins}
\end{equation}
Indeed a violating $h$ would give the nonzero lattice vector
$(h,h^Tw)$ of length at most $\sqrt5sB<4sB$. LLL and its reduced-basis
certificate have cost polynomial in $s+\log T+\log B$; there is no
enumeration of the tested polynomial family.

For the success analysis, enlarge the height to $R_s=2^s4sB$.
If every nonzero polynomial of height at most $R_s$ has value at least
$\eta$ in absolute value, then the test accepts whenever
\begin{equation}
 T\eta>(s+1)R_s.
 \label{eq:recourse-lattice-success}
\end{equation}
A lattice vector with $\norm{h}>R_s$ is already long. Otherwise its
last coordinate has magnitude at least
$T\eta-sR_s>R_s$. This proves the sufficient condition. Enlarging
the family is necessary because a short vector can have height greater
than $B$.

The coordinate polynomials $v_i$ and $1-v_i$ belong to the tested
family. Acceptance thus certifies that all free coordinates are truly
interior. With $\xi=sB/T$, set
\[
 \delta=\min\{1/2,\xi\},\qquad
 \mu=\min\{1,\xi/D\}.
\]
Every free coordinate is at least $\delta$ from its endpoints, and
every nonzero minor at $a$ has magnitude at least $\mu$.
Identically zero minor polynomials cause no difficulty. For a vertex
face, use $\delta=1/2$ and $\mu=1/D$ without LLL.

The
\href{../../research-20261002/new-direction/residual-convex-cubic-boundary.md}{conditional cubic fiber bound}
now computes a rational $\Gamma\ge1$, of bit length polynomial in
$I+\log T$, such that
\begin{equation}
 \dist(z,S_a)\le\Gamma
       [F(a,z)-\min_w F(a,w)]^{1/4}
 \quad\text{when the gap is in }[0,1].
 \label{eq:recourse-fiber-error}
\end{equation}
Its proof combines the cubic transverse estimate with a Hoffman bound
for \eqref{eq:recourse-fiber}. The supplied nonzero-minor margin
replaces the integer-minor bound for rational coefficients; the face
margin bounds the positive eigenvalues transverse to $\ker H_A$.
Rank-zero fibers are included. The lattice step supplies these premises
without knowing exact algebraic coordinates of $a$.

\subsection{Smoothed margins and the finite-grid transfer}

Choose a rational $U\ge1$ with $\nabla^2_{vv}F\preceq UI$ on the
box. Its numerical size will enter only the base precision, through
$\log U$. On any core face, the projected value
$V(v)=\min_zF(v,z)$ satisfies that $V(v)-U\norm{v}^2/2$ is concave.
Let $C$ be the interior points having an affine lower support on the
whole face. At $a\in C$ the support slope $g_a$ is unique and
\[
 0\le V(x)-V(a)-g_a^T(x-a)\le(U/2)\norm{x-a}^2.
\]
For $a,b\in C$, put $r=\norm{a-b}$. Comparing the upper support at
$a$ and lower support at $b$, then choosing
$h=r(g_b-g_a)/\norm{g_b-g_a}$ when this displacement stays in the
face, gives
\[
 (g_b-g_a)^Th\le(U/2)(r^2+\norm{h}^2),
 \qquad \norm{g_b-g_a}\le U\norm{a-b}.
\]
Thus the contact-gradient map is locally $U$-Lipschitz. A countable
interior cover and a disjoint measurable partition give
$|g(E\cap C)|\le U^m|E|$. An interior optimizer under tilt has this
contact slope, with the opposite sign. Continuous uniform noise
therefore puts probability at most $(U/(2\sigma))^m|E|$ on such cores.
No continuity of a residual optimizer is asserted.

For every nonzero integer polynomial $p$ of degree at most two,
\begin{equation}
 \bigl|\{v\in[0,1]^m:|p(v)|\le\eta\}\bigr|
 \le8\sqrt\eta,\qquad 0<\eta<1.
 \label{eq:recourse-sublevel}
\end{equation}
A nonzero square coefficient permits coordinate slicing, using the
univariate quadratic sublevel length bound $4\sqrt{\eta/|c|}$.
If only mixed quadratic terms occur, slice along
$(e_i+e_j)/\sqrt2$ for a nonzero mixed coefficient. The leading
magnitude is at least $1/2$ and the projected cube volume is
$\sqrt2$, giving the stated constant. Linear and constant polynomials
are easier.

Put $R=2^{s_*}4s_*B$ and
\[
 A_*=8\,3^k\max\{1,U/(2\sigma)\}^k(2R+1)^{s_*}.
\]
A union bound over faces and all degree-two integer polynomials of
height at most $R$ bounds the continuous bad-margin probability by
$A_*\sqrt\eta$. The logarithm of this family size is polynomial;
the algorithm never lists its members.

For a fixed face and polynomial, the bad event says that there exists
an interior face optimizer with $|p(a)|\le\eta$. It is expressed by
one existential full-variable witness block and one universal
competitor block, with degree at most three. The inherited
\href{../../research-20261002/new-direction/polynomial-finite-noise-tails.md}{finite-noise section bound}
therefore gives at most $2^{\poly(I)}$ interval or point components
in every one-dimensional noise section, uniformly in the fixed values
of the other coefficients. Replacing continuous marginals by grid
marginals one at a time changes each section probability by at most
twice its component count divided by $M$. Summing over the family gives
\begin{equation}
 \Pr_{\rm grid}\{\text{bad margin}\}
 \le A_*\sqrt\eta+C_*/M,
 \qquad C_*=2^{\poly(I)}.
 \label{eq:recourse-margin-tail}
\end{equation}
Exact-zero atoms are included. This argument bounds a sufficient
condition for LLL success; it does not need a semialgebraic description
of the LLL algorithm. The
\href{../cubic-recourse/margin-tail.md}{margin companion proof}
provides the complete section-count construction.

\subsection{One base certificate, one law, and a fixed-selector fallback}

Choose the exact full-point fallback factor $B_0\ge2$ first, then set
\[
 t=\min\{1/2,\sigma/(16kB_0)\},\qquad
 \eta=\min\{1/4,(16A_*B_0)^{-2}\}.
\]
Choose a power of two $T>(s_*+1)R/\eta$. All these quantities have
polynomial bit length and precede the noise mesh. Instantiate the
selected-core interface for $F$ and the $2k$ shifted baselines
$F\pm t v_i$. Their input lengths remain polynomial in $I$, and their
core second derivatives are unchanged. Use the same random vector
$\gamma$ for every instance. Choose one power-of-two $M$ satisfying
all these interfaces' lower mesh requirements and
\[
 (4k+C_*)/M\le1/(8B_0).
\]
Hence $\log M=\poly(I)$, independently of all later accuracy queries.
The shifted queries are auxiliary problems on the same draw, not new
random samples.

Run the face tests and one lattice test at the chosen scale. If the
lattice test fails, use exact fallback on the original draw. Equations
\eqref{eq:recourse-face-tail}, \eqref{eq:recourse-lattice-success},
and \eqref{eq:recourse-margin-tail} bound this failure probability by
$1/(4B_0)$. A missed endpoint forces rejection because its coordinate
polynomial vanishes. Every accepted certificate is valid independently
of the probability estimate.

The exact fallback selects the same $(a,p)$. A fixed-block formula
says that its witness is globally optimal and every equal-value
competitor has lexicographically larger core, or the same core and
residual squared norm at least that of the witness. Residual convexity
makes the witness unique. The degree remains at most three, and the
lexicographic comparison has linear-size Boolean syntax. The inherited
singleton-coordinate construction thus costs
$B_0\poly(I+b+q)$, with $B_0$ independent of $b,q$. Clipping the
rational approximation to the box preserves its distance guarantee.
Ties neither change the selector nor require resampling.

\subsection{Completing the residual point at every precision}

On acceptance, $\Gamma$ in \eqref{eq:recourse-fiber-error} has
polynomial base bit length. Given $\varepsilon=2^{-q}$, put
$e=\varepsilon/2$ and choose rational
$R_z\ge\max\{1,\sqrt n\}$ and
$G\ge\max\{1,\sup\norm{\nabla_vF}\}$. Set
\[
 \tau=\frac{e^6}{1024R_z^4\Gamma^4},\qquad
 \eta_q=\frac{\tau e^2}{8},\qquad
 d_q=\min\left\{\frac\varepsilon2,\frac{\tau e^2}{32G}\right\}.
\]
Obtain a short rational $b\in A$ with $\norm{b-a}\le d_q$, fixing
certified endpoints exactly, and solve the rational convex cubic
\[
 Q_b(z)=F(b,z)+\tau\norm z^2
\]
to objective gap at most $\eta_q$.

For the proof, let $z_\tau$ minimize
$F(a,z)+\tau\norm z^2$. Comparison with $p$ gives
$\norm{z_\tau}\le\norm p\le R_z$. If $s$ is its nearest point in
$S_a$ and $d=\norm{s-z_\tau}$, then the error bound and the
minimum-norm projection inequality imply
\[
 \frac{d^4}{\Gamma^4}
 \le F(a,z_\tau)-F(a,p)\le2R_z\tau d,
 \qquad \norm{z_\tau-p}^2\le2R_zd.
\]
The gap is initially at most $\tau R_z^2\le1$, so the error bound
applies. The chosen $\tau$ gives
$d\le e^2/(8R_z)$ and $\norm{z_\tau-p}\le e/2$.
Uniformly in $z$, $|F(b,z)-F(a,z)|\le Gd_q$. The returned point
therefore has true regularized gap at most
$\eta_q+2Gd_q\le\tau e^2/4$. Strong convexity places it within
$e/2$ of $z_\tau$, and hence within $e$ of $p$. Together with the
core error this proves the theorem's Euclidean bound. A finer point
query and the inherited value lower bound provide the optional sampled
objective-gap certificate.

All requested bit lengths are $\poly(I)+O(q)$. Add the common random
work factors of the original and $2k$ auxiliary core interfaces, and
add $B_0$ times the indicator of the single base certificate-failure
event. This sum controls every precision query and has expectation
\eqref{eq:recourse-work}. Independence between work factors and
certificate events is unnecessary. Large $U$, small certified margins,
and $\Gamma$ affect polynomial bit lengths rather than replacing $L$
by a larger numerical parameter in the expected-work factor.

The product box and total degree three are essential to this argument.
It covers, for example, $F(v,z)=vz^2$ and
$F(v,z)=v^2(2-z)$, which admit no finite quadratic core convexifier.
For the latter, the minimum-norm residual selection jumps from zero at
$v=0$ to one at every positive core. The proof avoids assuming that
selection is continuous. It does not extend the conclusion to arbitrary
coupled domains, quartics, expanded algebraic output, or the
unperturbed objective at the same cost.
